% thesis2paper deck. 10 slides for the assessment round, 5 for result rounds.
% One message per slide: the frame title IS the takeaway (a full sentence), the body is the evidence.
\documentclass[aspectratio=169,11pt]{beamer}
\usetheme[progressbar=frametitle,numbering=fraction]{metropolis}
\usepackage{booktabs,graphicx}
\newcommand{\fillme}[1]{\textcolor{red}{[#1]}}
\newcommand{\src}[1]{\par\vfill{\tiny\color{gray} Source: #1}}   % every result slide cites its source
\title{\fillme{Project}: \fillme{round / stage}}
\subtitle{\fillme{one-sentence headline}}
\date{\fillme{date}}

\begin{document}
\maketitle

% Suggested 10-slide arc for the assessment round:
%  1 The question and why it matters broadly
%  2 Setup / model / data in one picture (candidate for paper Fig. 1)
%  3-6 Key results, one per slide (figure + 1 line), each with \src{}
%  7 What reproduced and what did not
%  8 Positioning vs closest prior work
%  9 Gaps that block a top-journal paper
% 10 Proposed experiments / next steps + decisions needed
% Result rounds (5 slides): what we ran; 2-3 results; what changed in the ledger; next step.

\begin{frame}{\fillme{Takeaway sentence}}
  \begin{columns}[T]
    \column{0.58\textwidth}
      % \includegraphics[width=\linewidth]{fig_xxx.pdf}
      \fillme{figure}
    \column{0.40\textwidth}
      \begin{itemize}
        \item \fillme{evidence point}
      \end{itemize}
  \end{columns}
  \src{\fillme{raw/runs/E1/metrics.json, thesis p.42}}
\end{frame}

\end{document}
